\documentclass[12pt,a4paper]{article}

\usepackage[romanian]{babel}
\usepackage[T1]{fontenc}
\usepackage[utf8]{inputenc}
\usepackage{lmodern}
\usepackage{enumitem}

\title{Concepte Fundamentale în Bazele de Date}
\author{}
\date{}

\begin{document}

\maketitle

\section{RDBMS}
Un \textbf{RDBMS} (Relational Database Management System) este un sistem de
gestiune a bazelor de date care stochează informațiile sub formă de relații
(tabele). Exemple: MySQL, PostgreSQL, SQL Server, Oracle.

Caracteristici:
\begin{itemize}[noitemsep]
    \item model bazat pe tabele, chei primare și străine;
    \item integritate referențială;
    \item limbaj standardizat (SQL);
    \item tranzacții conforme cu modelul ACID.
\end{itemize}

\section{Modelul ACID}
Tranzacțiile dintr-un RDBMS trebuie să respecte proprietățile \textbf{ACID}:
\begin{itemize}[noitemsep]
    \item \textbf{Atomicitate}: tranzacția se execută complet sau deloc.
    \item \textbf{Consistență}: baza de date trece dintr-o stare validă în alta.
    \item \textbf{Izolare}: tranzacțiile nu se afectează reciproc.
    \item \textbf{Durabilitate}: modificările persistă după confirmare.
\end{itemize}

\section{SQL}
\textbf{SQL} (Structured Query Language) este limbajul standard pentru lucrul cu
baze de date relaționale. Permite:
\begin{itemize}[noitemsep]
    \item definirea structurii bazei de date;
    \item manipularea datelor;
    \item controlul accesului și al tranzacțiilor.
\end{itemize}

\section{NoSQL}
\textbf{NoSQL} reprezintă o categorie de sisteme de stocare non-relaționale,
optimizate pentru scalabilitate și flexibilitate. Tipuri:
\begin{itemize}[noitemsep]
    \item baze de date document (ex.: MongoDB);
    \item baze de date key--value (ex.: Redis);
    \item baze de date graf (ex.: Neo4j);
    \item baze de date column-store (ex.: Cassandra).
\end{itemize}

Caracteristici:
\begin{itemize}[noitemsep]
    \item schemă flexibilă;
    \item scalare orizontală;
    \item performanță ridicată pentru volume mari de date.
\end{itemize}

\section{LDD, LMD, LCD}
În SQL, operațiile sunt împărțite în trei categorii:

\subsection*{LDD -- Limbaj de Definire a Datelor}
Definește structura bazei de date.
\begin{itemize}[noitemsep]
    \item \texttt{CREATE}, \texttt{ALTER}, \texttt{DROP}
\end{itemize}

\subsection*{LMD -- Limbaj de Manipulare a Datelor}
Manipulează rândurile din tabele.
\begin{itemize}[noitemsep]
    \item \texttt{SELECT}, \texttt{INSERT}, \texttt{UPDATE}, \texttt{DELETE}
\end{itemize}

\subsection*{LCD -- Limbaj de Control al Datelor}
Controlează accesul și permisiunile.
\begin{itemize}[noitemsep]
    \item \texttt{GRANT}, \texttt{REVOKE}
\end{itemize}

\section{Designul Bazelor de Date}
Designul unei baze de date urmărește:
\begin{itemize}[noitemsep]
    \item identificarea entităților și atributelor;
    \item definirea relațiilor și a cheilor;
    \item normalizarea (1NF, 2NF, 3NF etc.);
    \item optimizarea performanței și a integrității;
    \item proiectarea fizică (indexuri, partiționare).
\end{itemize}

Procesul tipic:
\begin{enumerate}[noitemsep]
    \item analiza cerințelor;
    \item model conceptual (ERD);
    \item model logic (relații, chei);
    \item model fizic (implementare în RDBMS).
\end{enumerate}

\end{document}
